% Options for packages loaded elsewhere
\PassOptionsToPackage{unicode}{hyperref}
\PassOptionsToPackage{hyphens}{url}
\PassOptionsToPackage{dvipsnames,svgnames,x11names}{xcolor}
%
\documentclass[
  letterpaper,
  DIV=11,
  numbers=noendperiod]{scrartcl}

\usepackage{amsmath,amssymb}
\usepackage{iftex}
\ifPDFTeX
  \usepackage[T1]{fontenc}
  \usepackage[utf8]{inputenc}
  \usepackage{textcomp} % provide euro and other symbols
\else % if luatex or xetex
  \usepackage{unicode-math}
  \defaultfontfeatures{Scale=MatchLowercase}
  \defaultfontfeatures[\rmfamily]{Ligatures=TeX,Scale=1}
\fi
\usepackage{lmodern}
\ifPDFTeX\else  
    % xetex/luatex font selection
\fi
% Use upquote if available, for straight quotes in verbatim environments
\IfFileExists{upquote.sty}{\usepackage{upquote}}{}
\IfFileExists{microtype.sty}{% use microtype if available
  \usepackage[]{microtype}
  \UseMicrotypeSet[protrusion]{basicmath} % disable protrusion for tt fonts
}{}
\makeatletter
\@ifundefined{KOMAClassName}{% if non-KOMA class
  \IfFileExists{parskip.sty}{%
    \usepackage{parskip}
  }{% else
    \setlength{\parindent}{0pt}
    \setlength{\parskip}{6pt plus 2pt minus 1pt}}
}{% if KOMA class
  \KOMAoptions{parskip=half}}
\makeatother
\usepackage{xcolor}
\setlength{\emergencystretch}{3em} % prevent overfull lines
\setcounter{secnumdepth}{5}
% Make \paragraph and \subparagraph free-standing
\ifx\paragraph\undefined\else
  \let\oldparagraph\paragraph
  \renewcommand{\paragraph}[1]{\oldparagraph{#1}\mbox{}}
\fi
\ifx\subparagraph\undefined\else
  \let\oldsubparagraph\subparagraph
  \renewcommand{\subparagraph}[1]{\oldsubparagraph{#1}\mbox{}}
\fi


\providecommand{\tightlist}{%
  \setlength{\itemsep}{0pt}\setlength{\parskip}{0pt}}\usepackage{longtable,booktabs,array}
\usepackage{calc} % for calculating minipage widths
% Correct order of tables after \paragraph or \subparagraph
\usepackage{etoolbox}
\makeatletter
\patchcmd\longtable{\par}{\if@noskipsec\mbox{}\fi\par}{}{}
\makeatother
% Allow footnotes in longtable head/foot
\IfFileExists{footnotehyper.sty}{\usepackage{footnotehyper}}{\usepackage{footnote}}
\makesavenoteenv{longtable}
\usepackage{graphicx}
\makeatletter
\def\maxwidth{\ifdim\Gin@nat@width>\linewidth\linewidth\else\Gin@nat@width\fi}
\def\maxheight{\ifdim\Gin@nat@height>\textheight\textheight\else\Gin@nat@height\fi}
\makeatother
% Scale images if necessary, so that they will not overflow the page
% margins by default, and it is still possible to overwrite the defaults
% using explicit options in \includegraphics[width, height, ...]{}
\setkeys{Gin}{width=\maxwidth,height=\maxheight,keepaspectratio}
% Set default figure placement to htbp
\makeatletter
\def\fps@figure{htbp}
\makeatother

\KOMAoption{captions}{tableheading}
\makeatletter
\makeatother
\makeatletter
\makeatother
\makeatletter
\@ifpackageloaded{caption}{}{\usepackage{caption}}
\AtBeginDocument{%
\ifdefined\contentsname
  \renewcommand*\contentsname{Table of contents}
\else
  \newcommand\contentsname{Table of contents}
\fi
\ifdefined\listfigurename
  \renewcommand*\listfigurename{List of Figures}
\else
  \newcommand\listfigurename{List of Figures}
\fi
\ifdefined\listtablename
  \renewcommand*\listtablename{List of Tables}
\else
  \newcommand\listtablename{List of Tables}
\fi
\ifdefined\figurename
  \renewcommand*\figurename{Figure}
\else
  \newcommand\figurename{Figure}
\fi
\ifdefined\tablename
  \renewcommand*\tablename{Table}
\else
  \newcommand\tablename{Table}
\fi
}
\@ifpackageloaded{float}{}{\usepackage{float}}
\floatstyle{ruled}
\@ifundefined{c@chapter}{\newfloat{codelisting}{h}{lop}}{\newfloat{codelisting}{h}{lop}[chapter]}
\floatname{codelisting}{Listing}
\newcommand*\listoflistings{\listof{codelisting}{List of Listings}}
\makeatother
\makeatletter
\@ifpackageloaded{caption}{}{\usepackage{caption}}
\@ifpackageloaded{subcaption}{}{\usepackage{subcaption}}
\makeatother
\makeatletter
\@ifpackageloaded{tcolorbox}{}{\usepackage[skins,breakable]{tcolorbox}}
\makeatother
\makeatletter
\@ifundefined{shadecolor}{\definecolor{shadecolor}{rgb}{.97, .97, .97}}
\makeatother
\makeatletter
\makeatother
\makeatletter
\makeatother
\ifLuaTeX
  \usepackage{selnolig}  % disable illegal ligatures
\fi
\IfFileExists{bookmark.sty}{\usepackage{bookmark}}{\usepackage{hyperref}}
\IfFileExists{xurl.sty}{\usepackage{xurl}}{} % add URL line breaks if available
\urlstyle{same} % disable monospaced font for URLs
\hypersetup{
  pdftitle={MKTG100 Numeracy},
  pdfauthor={AD},
  colorlinks=true,
  linkcolor={blue},
  filecolor={Maroon},
  citecolor={Blue},
  urlcolor={Blue},
  pdfcreator={LaTeX via pandoc}}

\title{MKTG100 Numeracy}
\usepackage{etoolbox}
\makeatletter
\providecommand{\subtitle}[1]{% add subtitle to \maketitle
  \apptocmd{\@title}{\par {\large #1 \par}}{}{}
}
\makeatother
\subtitle{Solutions To Week 7b's Exercises}
\author{AD}
\date{}

\begin{document}
\maketitle
\ifdefined\Shaded\renewenvironment{Shaded}{\begin{tcolorbox}[interior hidden, borderline west={3pt}{0pt}{shadecolor}, enhanced, boxrule=0pt, breakable, frame hidden, sharp corners]}{\end{tcolorbox}}\fi

\renewcommand*\contentsname{Table of contents}
{
\hypersetup{linkcolor=}
\setcounter{tocdepth}{3}
\tableofcontents
}
\hypertarget{week-7b-exponents-and-roots}{%
\section{Week 7b: Exponents and
Roots}\label{week-7b-exponents-and-roots}}

\hypertarget{learning-goals}{%
\subsection{Learning Goals}\label{learning-goals}}

\begin{itemize}
\tightlist
\item
  Know how to evaluate exponential expressions
\item
  Know how to write exponential expressions when multiplying or dividing
  with the same base.
\item
  Know how to find roots.
\end{itemize}

\hypertarget{exponents}{%
\subsection{Exponents}\label{exponents}}

Rules:

\begin{enumerate}
\def\labelenumi{\arabic{enumi}.}
\item
  \(a^m \times a^n = a^{m+n}\)
\item
  \(\frac{a^m}{a^n} = a^{m-n}\)
\item
  \((a^m)^n=a^{mn}\)
\item
  \(a^{-m}=\frac{1}{a^{m}}\)
\end{enumerate}

\hypertarget{problems}{%
\subsubsection{Problems}\label{problems}}

Simply:

\begin{enumerate}
\def\labelenumi{\arabic{enumi}.}
\item
  \(x^2 \times x^3\) = \ldots{}
\item
  \(\frac{3^5}{3^2}\) = \ldots{}
\item
  \((4^5)^2\) = \ldots{}
\item
  \(2^{-3}\) = \ldots{}
\end{enumerate}

\hypertarget{watch}{%
\subsubsection{Watch}\label{watch}}

\hypertarget{exercises}{%
\subsubsection{Exercises}\label{exercises}}

Simply:

\begin{enumerate}
\def\labelenumi{\arabic{enumi}.}
\item
  \(z^3 \times z^2\) = \ldots{}
\item
  \(\frac{4^3}{4^2}\) = \ldots{}
\item
  \((5^2)^3\) = \ldots{}
\item
  \(3^{-2}\) = \ldots{}
\end{enumerate}

\hypertarget{attempt-the-exercises-on-moodle}{%
\subsection{Attempt the Exercises on
Moodle}\label{attempt-the-exercises-on-moodle}}

At this point, you would have worked on the above exercise questions.

Now, I want you to attempt these questions on Moodle.

Please click the link below:

\href{https://modules.lancaster.ac.uk/mod/quiz/view.php?id=2012499}{Week
7a Exercises on Moodle}

Once you submit all your answers on Moodle, you complete your learning
on this week's topic.

\hypertarget{roots}{%
\subsection{Roots}\label{roots}}

\textbf{Square roots}. 4 is called a square root of 16 because
\(4^2=16\). The symbol \(\sqrt[2]{}\) or \(\sqrt{}\) is used to
represent square root. For example, \(\sqrt{16}=4\). \(\sqrt{16}\) can
be written as \(16^\frac{1}{2}\).

\textbf{Cube roots}. 2 is called a cube root of 8 because \(2^3=8\). The
symbol \(\sqrt[3]{}\) is used to represent cube root. For example,
\(\sqrt[3]{8}=2\). \(\sqrt[3]{8}\) can be written as \(8^\frac{1}{3}\).

** In general, \(\sqrt[n]{a}\) can be written as \(a^\frac{1}{n}\). It
is hard sometimes to evaluate this expression manually. You can use a
scientific calculator or \texttt{Excel} to do it. For example, by typing
\texttt{=16\^{}(1/4)} in \texttt{Excel}, it gives 2. That is
\(\sqrt[4]{16}=16^\frac{1}{4}=2\). You can prove that
\(2 \times 2 \times 2 \times 2 = 16\)

\hypertarget{problems-1}{%
\subsubsection{Problems}\label{problems-1}}

Evaluate \(\sqrt[3]{16}\) (you can use a scientific calculator or
\texttt{Excel})

\hypertarget{watch-1}{%
\subsubsection{Watch}\label{watch-1}}

\hypertarget{exercises-1}{%
\subsubsection{Exercises}\label{exercises-1}}

Evaluate

\begin{enumerate}
\def\labelenumi{\arabic{enumi}.}
\item
  \(\sqrt[3]{81}\)
\item
  \(\sqrt[4]{256}\)
\item
  \(\sqrt{81}\)
\item
  \(\sqrt[3]{216}\)
\item
  \(\sqrt[4]{625}\)
\item
  \(\sqrt[3]{15.63}\)
\end{enumerate}

\hypertarget{attempt-the-exercises-on-moodle-1}{%
\subsection{Attempt the Exercises on
Moodle}\label{attempt-the-exercises-on-moodle-1}}

At this point, you would have worked on the above exercise questions.

Now, I want you to attempt these questions on Moodle.

Please click the link below:

\href{https://modules.lancaster.ac.uk/mod/quiz/view.php?id=2235175}{Week
7b Exercises on Moodle}

Once you submit all your answers on Moodle, you complete your learning
on this week's topic.



\end{document}
